\section{Discussion}
\label{sec:discussion}

We first interpret the answers to the research questions. We then discuss when whole-paper prompting is enough, how reliable the extracted gaps are, why section rules do not transfer to computer-science papers, and what the results imply for evaluation data. The section ends with the limitations of the study.

\subsection{Answers to the research questions}

\paragraph{RQ1: each stage has its own role.} Stage~A is cheap, but it is useful only where section headings are informative. Stage~B sets the upper limit on recall, because Stage~C only removes candidates and cannot recover a gap that Stage~B missed. Stage~C provides most of the gain in precision. Throughout, every reported gap remains a sentence the authors wrote.

\paragraph{RQ1: a training-free Stage~B is viable but less precise.} A zero-shot NLI classifier needs no gap-labelled data and works on both biomedical and computer-science text. On SAL, a supervised classifier was 6.9 to 9.0 percentage points more precise in every mode, while the $F_1$ difference in high-precision mode was small and not significant. This supervised comparison is available only for SAL, and its classifier was selected on the SAL test split. Fine-tuning on LLM labels did not substitute for human labels.

\paragraph{RQ1: one instruction sets the trade-off.} The three modes kept their order on every evaluation set, on the held-out SAL dev papers and on the gold validation split before any test run. Within these datasets and prompts, the mode paragraph therefore provides a practical control. High precision suits idea generation, where every false gap costs a downstream model or a person time, and high recall suits a reviewer who wants to see every candidate. High precision costs recall mostly on research gaps. A possible reason, which we have not tested, is that the high-precision instruction asks the model to drop background and motivation, and authors often state open problems in exactly that role.

\paragraph{RQ1: the combination matters most on full papers.} Compared with published sentence-level methods, the funnel was far ahead on SAL-40 and ahead on arXiv, where gap sentences are rare and spread through long papers. On FWS, which consists of conclusion sections and offers plentiful training data, a TF-IDF classifier with logistic regression came close to the funnel. The funnel is therefore most useful on full papers. On short, pre-selected text with plentiful training data, a trained classifier is a serious alternative.

\paragraph{RQ2: the funnel helps the 9B and 4B models on long papers.} With both Qwen3.5-9B and Qwen3.5-4B, whole-paper prompting performed poorly on long arXiv papers, and the funnel did not. With the 27B model the picture is mixed. The funnel scored higher on the two human-labelled benchmarks, using about one-third of the input tokens, and lower on arXiv papers against LLM labels. The caveats of Section~\ref{sec:comparison} limit what the 27B results show, and the SAL differences were not significant. The call format matters for cost as well: sending each paper's candidates in a single call cuts the input on full papers to about one-seventh of whole-paper prompting on arXiv and to one-eighth on SAL-40.

\paragraph{RQ3: LLM labels are not yet a substitute for human labels.} On SAL, the LLM labels include nearly every human-labelled limitation, but they also mark many sentences the annotators left unlabelled. The same output therefore scores 66.2\% precision against human labels and 93.2\% against LLM labels, and a classifier fine-tuned on LLM labels improves against those labels but not against human ones. Within these data, precision measured against LLM labels should not be read as precision against human judgement, which is why we call the three-type results provisional. Whether the same holds for future work and research gaps is what the planned human study will test.

\subsection{When is whole-paper prompting enough?}
\label{sec:whole-discussion}

Three factors decide the comparison in our data.
\begin{itemize}
  \item \textbf{Input length.} SAL papers have about 200 sentences and FWS excerpts about seven. arXiv papers have a median of 239 sentences and up to 1526, with whole-paper inputs of up to 72,000 tokens. The 9B and 4B models handled the short inputs but selected far too much on several long arXiv papers. In the median-split analysis, which was run for the 27B and 4B models, the funnel helped the 4B model on the longer half and not on the shorter half.
  \item \textbf{Model size and generation.} The 27B model handled every input, whereas the 9B and 4B models both failed on long inputs in a similar way. These two models belong to the same generation, Qwen3.5, so the failure is unlikely to be specific to the 4B checkpoint. The 27B model, however, belongs to the newer Qwen3.8 generation. Size is therefore not the only difference between it and the two smaller models, and our results cannot separate the effect of size from the effect of model design.
  \item \textbf{Reference labels.} The arXiv labels were produced by LLMs that read each whole paper in one call, the same input format as the whole-paper baseline, and may therefore favour it. On the two benchmarks labelled by people, the funnel was ahead with the 27B model.
\end{itemize}

Taken together, these factors give a simple answer for our data. With the 27B model, whole-paper prompting was a reasonable choice, especially on shorter papers, as long as its higher input cost is acceptable. With the 9B and 4B models, it held up on SAL-40 but failed on several long arXiv papers, and there these models needed the funnel.

The call format also matters. The funnel's main cost on arXiv is that each candidate carries its own context and receives its own verdict. Sending each paper's candidates in a single call removes most of this cost. This one-call-per-paper format is strong with the 9B model on all three datasets and is the most accurate setup on SAL-40 with the 4B model. With the 27B model, however, it loses up to 0.05 $F_1$ on arXiv, and asking for a list of kept sentences instead of a verdict on each candidate makes the model less selective. Cheaper formats are thus a trade-off, not a free improvement.

\subsection{How reliable are the extracted gaps?}
\label{sec:quality}

How good the extracted gaps are depends on which labels are trusted. On SAL-40, high-precision mode reaches 66.2\% precision against SAL's labels and 93.2\% against the LLM labels. Either label set may be wrong on the 77 sentences where they disagree. On the one hand, SAL's annotators may have left some limitation statements unlabelled. On the other hand, the LLM labels are inclusive by construction, because adjudication added sentences that only one model had marked. The two values describe the same output under two references, and neither is a bound on its true precision.

Label errors may also affect the gold corpus. In an earlier high-recall run during development, which is not part of the results above, five of the 34 test papers accounted for 59\% of the false positives. Each of these papers had only about 18--24 gold gaps despite having 220 to 824 sentences, which may indicate that long papers are under-labelled. If so, some sentences counted as false positives are in fact gaps, and precision against these labels is understated, most of all for systems that keep many sentences. The effect on measured recall is less clear, because unlabelled gaps are missing from both the numerator and the denominator.

% HUMAN STUDY: replace the next sentence by the discussion text in docs/publication/HUMAN_STUDY_INSERT.md.
The planned human study (Section~\ref{sec:methods}) will measure precision against human judgement on a frozen sample from arXiv and SAL-40.

\subsection{Section rules do not transfer to computer-science papers}
\label{sec:localisation}

In RCT reports, limitations cluster in the Discussion and Conclusion sections, so a section rule works. In arXiv papers it did not work in our development experiments. An earlier localiser kept only the sections where gaps are usually stated, such as the discussion and the conclusion. We measured it against the limitation and future-work labels of the arXiv papers, with a fuzzy sentence match. It kept only 55.6\% of these gaps when sections were inferred from the extracted text, and 55.2\% when sections were taken from the LaTeX sources of 175 papers. Better section information therefore did not help: with LaTeX sections, 82\% of the missed gaps were in methods, experiments or results sections. This development analysis used a different target and matching rule from our main evaluation and was not repeated for this version. For arXiv papers we therefore use only the junk filter as Stage~A, which keeps 97\% of the gold test gaps, and leave the narrowing to Stage~B.

\subsection{Evaluation data needs checking}

The FWS case artefact shows that a published benchmark can contain a trivial signal that a modern model will find. FWS results from case-sensitive models, or from LLMs given the raw text, should be treated with caution unless the text is normalised. More generally, the precision we measure on SAL-40 changed more with the label source than with the choice of Stage~B. Stating which labels a number rests on is as important as reporting the number.

\subsection{Limitations}

\begin{itemize}
  \item \textbf{LLM labels.} All results on the gold corpus rest on labels produced by Claude models, which have not yet been checked by people. The agreement between the two labelling models ($\kappa = 0.778$) measures their consistency, not their accuracy. Claude models also assisted in writing the prompts and this manuscript, so the labels and the prompts may share biases.
  \item \textbf{Development on test papers.} The Stage~C model, batch size, use of context and high-recall instruction were chosen on 196 candidates from 27 of the 34 gold test papers, and the balanced and high-precision instructions were written after reading false positives on gold test papers. We report results without these candidates and on the 7 untouched papers, ran and report all modes on every dataset, and checked the modes on the validation split first. Still, the gold test split is not untouched. SAL dev provides the cleanest held-out check, and its results agree with those on SAL-40.
  \item \textbf{Fine-tuned classifier and arXiv scores.} The one-call-per-paper funnel uses a Stage~B classifier fine-tuned on the gold training labels, which may favour it when it is scored against the same kind of labels on arXiv. The zero-shot funnel with 20 candidates per call shows a similar gain with the 4B model ($F_1$ 0.544 against 0.296 for whole-paper prompting), so the gain does not depend on fine-tuning.
  \item \textbf{Comparisons designed later.} The comparisons with whole-paper prompting, smaller models, other formats and published methods were designed after the main results were known. The whole-paper baseline and the published-method baselines were frozen before their test runs, and all runs of these comparisons are reported. The comparison with whole-paper prompting is also subject to the three caveats of Section~\ref{sec:comparison}.
  \item \textbf{Published methods.} The rules after Kilicoglu et al. and the classifier after Zhang et al. are adaptations, not exact reproductions. We could not reproduce the detector of Lan et al., and our fine-tuned SciBERT and BiomedBERT runs on arXiv failed, so these are not among the baselines.
  \item \textbf{Statistics.} Apart from two stability checks (one repeated high-recall run and 70 repeated batches), each configuration was run once, and the many paired comparisons are not corrected for multiplicity. Small differences, especially those whose confidence intervals include zero, should not be treated as established effects.
  \item \textbf{Scope.} SAL covers RCT reports only, FWS covers NLP conclusion sections only, and each labels a single gap type. The six biomedical papers in the gold test split are also SAL papers. The arXiv papers are a convenience sample of 28 test papers. Stage~C was tested with one model family, in which the largest model belongs to a newer generation than the two smaller ones.
  \item \textbf{Research gaps.} Half of the gold labels are research gaps, which the earlier stages retain better than the other types. We therefore also report recall by type and scores without research gaps.
  \item \textbf{Cost.} Cost is measured as LLM tokens and call time on our hardware, not in money, and excludes Stages~A and~B.
\end{itemize}
